\begin{abstract}
\ifmetaum
Este relatório descreve a primeira fase de um projeto de Processamento de Linguagem Natural, centrado no domínio da Cibersegurança em Língua Portuguesa. Nesta etapa, procedeu-se à construção de um conjunto de dados composto por pares de pergunta-resposta curta, abrangendo tanto conceitos tradicionais de cibersegurança como aspetos emergentes de segurança em sistemas de Inteligência Artificial, nomeadamente vulnerabilidades associadas a Large Language Models (LLMs). A criação do corpus combinou a formulação manual de questões por parte dos elementos do grupo com a geração automática através do modelo AMALIA (q4\_K\_M)~\cite{amalia_ollama}, executado localmente via Ollama. As questões geradas automaticamente foram sujeitas a uma curadoria rigorosa, assegurando critérios de clareza, ausência de redundância e validade factual das respostas. O conjunto de dados final, disponibilizado em formato .tsv, contém a identificação do tema, a pergunta, o método de geração e a respetiva resposta, constituindo a base empírica para as etapas seguintes do projeto.
\fi
\end{abstract}
